В рамках выпускной квалификационной работы была решена задача программной реализации высокопроизводительных алгоритмов эллиптической криптографии для управляемой среды выполнения. Разработана криптографическая библиотека, реализующая примитивы Curve25519 и Ed25519 для платформы .NET. Основными результатами работы являются:

\begin{enumerate}
    \item Реализация математического ядра на базе Radix-51. Создан модуль арифметики для конечного поля $GF(2^{255}-19)$. Использование 64-битных регистров с отложенным переносом позволило сократить общее количество базовых вычислительных операций по сравнению с классическими реализациями.
    
    \item Обеспечение константного времени выполнения. Протоколы обмена ключами и цифровой подписи (согласно RFC 7748 и RFC 8032) реализованы с использованием проективных координат и полных законов сложения. Исключение условных переходов защищает библиотеку от атак по времени (timing attacks).
    
    \item Исключение аллокаций памяти. Применение типов ref struct и Span<byte> позволило реализовать все криптографические вычисления без выделения памяти в управляемой куче. Это полностью снимает нагрузку со сборщика мусора (Garbage Collector) во время работы алгоритма.
    
    \item Внедрение пакетной обработки (SIMD). Архитектура данных была перепроектирована под формат SoA (Structure of Arrays). С помощью интринсиков System.Runtime.Intrinsics.X86.Avx2 реализовано параллельное выполнение четырех независимых операций скалярного умножения.
    
    \item Оценка производительности. Бенчмаркинг подтвердил работоспособность предложенной архитектуры: разработанная библиотека опережает индустриальный стандарт Bouncy Castle по пропускной способности в среднем на 30–40\% при использовании пакетной обработки.
\end{enumerate}

Цель работы — проектирование и реализация быстрой, безопасной и кроссплатформенной криптографической библиотеки X25519/Ed25519 для .NET — достигнута. Поставленные задачи выполнены в полном объеме.

Практическая значимость работы заключается в следующем:
\begin{enumerate}
    \item Создана готовая к использованию библиотека на чистом C\# (managed code), которую можно интегрировать в высоконагруженные серверные приложения без привязки к нативным C/C++ зависимостям (P/Invoke).
    \item Показано, что применение интринсиков и ручное управление памятью на стеке позволяет управляемому коду достигать производительности, достаточной для конкуренции с системными криптографическими библиотеками.
    \item Предложенная архитектура векторизации (SoA) может быть переиспользована для ускорения других алгоритмов асимметричной криптографии в экосистеме .NET.
\end{enumerate}

Основные выводы:
\begin{enumerate}
    \item Платформа .NET 10 предоставляет необходимый низкоуровневый инструментарий для создания эффективных криптографических примитивов без потери производительности.
    \item Переход к пакетной обработке с помощью SIMD-инструкций дает существенный прирост скорости, однако требует изменения подхода к хранению структур данных в памяти.
    \item Отказ от аллокаций (Zero Allocation) на критическом пути выполнения является обязательным условием для обеспечения стабильной пропускной способности и предсказуемого времени отклика сервера.
\end{enumerate}

Перспективы развития проекта включают расширение поддержки аппаратных архитектур (инструкции AVX-512 и ARM NEON), добавление поддержки криптографической группы Ristretto255, а также реализацию адаптеров для совместимости со стандартными интерфейсами пространства имен System.Security.Cryptography.